

\section{One-Sided Limits}
\label{sec:1.2}

When we compute the limit of a function \(f\) as \(x\) approaches \(a\), we observe its behavior from both sides. However, a function may behave differently from the left than from the right, as happens with piecewise functions. A function may also be undefined on one side of \(a\), leaving only one direction to observe. These situations lead to the concept of a one-sided limit. \\
 
By the end of this section, the student will be able to:
\begin{itemize}
    \item interpret the one-sided limit of a function through graphs and tables of values;
    \item evaluate one-sided limits of functions; and
    \item determine the limit of piecewise functions using one-sided limits.
\end{itemize}

\medskip
We first illustrate the idea with two concrete functions.

\medskip
\textbf{Illustration 1.2.1.}
Let
\[
f(x) = \begin{cases}
3 - 5x^{2}, & x < 1, \\[4pt]
4x - 3,     & x \ge 1.
\end{cases}
\]
If \(x\) is less than 1 but very close to 1, then \(f(x) = 3 - 5x^{2}\) and its values approach \(-2\). If \(x\) approaches 1 through values greater than 1, we use \(f(x) = 4x - 3\) and \(f(x)\) approaches \(1\). The two sides do not agree, so \(\lim_{x \to 1} f(x)\) does not exist.

\begin{center}
\begin{tikzpicture}[scale=0.85]
    \draw[->,thick] (-2.5,0) -- (3.2,0) node[below] {\(x\)};
    \draw[->,thick] (0,-3.2) -- (0,3.8) node[left] {\(y\)};
    \foreach \x in {-2,-1,1,2}
        \draw (\x,0.1) -- (\x,-0.1) node[below] {\(\x\)};
    \foreach \y in {-2,-1,1,2,3}
        \draw (0.1,\y) -- (-0.1,\y) node[left] {\(\y\)};
    \node at (0,0) [below left] {\(0\)};
    % y = 3 - 5x^2, x < 1
    \draw[very thick, primaryColor, domain=-1.05:0.97, smooth]
        plot (\x, {3 - 5*\x*\x});
    \draw[primaryColor, fill=white, thick] (1,-2) circle (2.5pt);
    %y = 4x - 3, x >= 1
    \draw[very thick, primaryColor, domain=1:1.65, smooth]
        plot (\x, {4*\x - 3});
    \fill[primaryColor] (1,1) circle (2.5pt);
    %  x=1
    \draw[dashed, gray] (1,-2.8) -- (1,3.2);

    \node[primaryColor, anchor=west] at (1.5,2.8) {\(y = 3 - 5x^{2}\)};
    \node[primaryColor, anchor=south west] at (1.05,1.2) {\(y = 4x - 3\)};
\end{tikzpicture}

\smallskip
\textbf{Figure 2:} Graph of \(y = f(x)\) in Illustration~1.2.1.
\end{center}

\medskip
\textbf{Illustration 1.2.2.}
Let \(f(x) = \sqrt{x}\). The domain is \([0, \infty)\), so we can only consider values of \(x\) greater than 0. We cannot observe the behavior from the left of 0 at all.

\begin{center}
\begin{tikzpicture}[scale=0.85]
    \draw[->,thick] (-0.5,0) -- (5.2,0) node[below] {\(x\)};
    \draw[->,thick] (0,-0.5) -- (0,3.2) node[left] {\(y\)};
    \foreach \x in {1,2,3,4}
        \draw (\x,0.1) -- (\x,-0.1) node[below] {\(\x\)};
    \foreach \y in {1,2}
        \draw (0.1,\y) -- (-0.1,\y) node[left] {\(\y\)};
    \node at (0,0) [below left] {\(0\)};
    \draw[very thick, primaryColor, domain=0:4.5, smooth]
        plot (\x, {sqrt(\x)});
    \fill[primaryColor] (0,0) circle (2.5pt);
    \node[primaryColor, anchor=south west] at (3.2,2.2) {\(y = \sqrt{x}\)};
\end{tikzpicture}

\smallskip
\textbf{Figure 3:} Graph of \(y = \sqrt{x}\) in Illustration~1.2.2.
\end{center}

\bigskip

These cases lead us to the formal definition of one-sided limits.

\medskip
Let \(f\) be a function defined at every number in some open interval \((c, a)\). The \textbf{limit of \(f(x)\) as \(x\) approaches \(a\) from the left} is \(L\), denoted
\[
\lim_{x \to a^{-}} f(x) = L,
\]
if the values of \(f(x)\) get closer and closer to \(L\) as \(x\) gets closer and closer to \(a\), but with \(x\) restricted to values less than \(a\).

Similarly, let \(f\) be defined on some open interval \((a, c)\). The \textbf{limit of \(f(x)\) as \(x\) approaches \(a\) from the right} is \(L\), denoted
\[
\lim_{x \to a^{+}} f(x) = L,
\]
if the values of \(f(x)\) get closer and closer to \(L\) as \(x\) gets closer and closer to \(a\), but with \(x\) restricted to values greater than \(a\).


\begin{remarkbox}{Remark 1}

\begin{enumerate}
    \item The basic limit laws in Section~1.1 remain valid when \(x \to a\) is replaced by \(x \to a^{-}\) or \(x \to a^{+}\).
    \item The expression \(\displaystyle\lim_{x \to a} f(x)\) is sometimes called the \textbf{two-sided limit} to distinguish it from one-sided limits.
\end{enumerate}
\end{remarkbox}


\begin{examplebox}{Example 2}

Let \(H(x)\) be the Heaviside step function:
\[
H(x) = \begin{cases}
1, & x \ge 0, \\[4pt]
0, & x < 0.
\end{cases}
\]
For \(x > 0\), \(H(x) = 1\) identically, so \(\displaystyle\lim_{x \to 0^{+}} H(x) = \lim_{x \to 0^{+}} 1 = 1\). For \(x < 0\), \(H(x) = 0\) identically, so \(\displaystyle\lim_{x \to 0^{-}} H(x) = \lim_{x \to 0^{-}} 0 = 0\).
\end{examplebox}


\begin{examplebox}{Example 3}

Return to the function from Illustration~1.2.1:
\[
f(x) = \begin{cases}
3 - 5x^{2}, & x < 1, \\[4pt]
4x - 3,     & x \ge 1.
\end{cases}
\]
\[
\lim_{x \to 1^{-}} f(x) = \lim_{x \to 1^{-}} (3 - 5x^{2}) = -2,
\qquad
\lim_{x \to 1^{+}} f(x) = \lim_{x \to 1^{+}} (4x - 3) = 1.
\]
\end{examplebox}

\bigskip


\subsection*{Approaching Zero Through Positive or Negative Values}
\label{subsec:1.2.1}


Suppose a function \(f\) satisfies \(\lim_{x \to a} f(x) = 0\). We can distinguish the manner in which \(f\) approaches 0. Consider \(f(x) = 2 - x\). As \(x \to 2^{-}\), we have \(f(x) \to 0\) but passing through \emph{positive} values. As \(x \to 2^{+}\), \(f(x)\) approaches 0 through \emph{negative} values:

\begin{center}
\small
\begin{tabular}{c|cccc|c|cccc}
\(x\)      & \(1.9\) & \(1.99\) & \(1.999\) & \(1.9999\) & \(2\) & \(2.0001\) & \(2.001\) & \(2.01\) & \(2.1\) \\
\hline
\(f(x)\)   & \(0.1\) & \(0.01\) & \(0.001\) & \(0.0001\) & \(0\) & \(-0.0001\) & \(-0.001\) & \(-0.01\) & \(-0.1\)
\end{tabular}
\normalsize
\end{center}

Now consider \(g(x) = (2 - x)^{2}\). Again \(\lim_{x \to 2} g(x) = 0\), but this time \(g(x)\) passes through \emph{positive} values from \emph{both} sides:

\begin{center}
\small
\begin{tabular}{c|cccc|c|cccc}
\(x\)      & \(1.9\)   & \(1.99\)     & \(1.999\)       & \(1.9999\)       & \(2\) & \(2.0001\)       & \(2.001\)       & \(2.01\)     & \(2.1\)   \\
\hline
\(g(x)\)   & \(0.01\)  & \(0.0001\)   & \(0.000001\)    & \(0.00000001\)   & \(0\) & \(0.00000001\)   & \(0.000001\)    & \(0.0001\)   & \(0.01\)
\end{tabular}
\normalsize
\end{center}

This distinction matters when we take even roots of quantities that tend to zero.


\begin{definitionbox}{Definition 4}

Suppose \(\displaystyle\lim_{x \to a} f(x) = 0\). If \(f(x)\) approaches 0 through positive values, we write
\[
f(x) \to 0^{+}.
\]
If \(f(x)\) approaches 0 through negative values, we write
\[
f(x) \to 0^{-}.
\]
\end{definitionbox}


\begin{remarkbox}{Remark 5}

\begin{itemize}
    \item If \(f(x) \to 0^{+}\) as \(x \to a\) and \(n\) is even, then \(\displaystyle\lim_{x \to a} \sqrt[n]{f(x)} = 0\).
    \item If \(f(x) \to 0^{-}\) as \(x \to a\) and \(n\) is even, then \(\displaystyle\lim_{x \to a} \sqrt[n]{f(x)}\) does not exist.
\end{itemize}
\end{remarkbox}


\begin{examplebox}{Example 6}

Let \(f(x) = \sqrt{2 - x}\). As \(x \to 2^{-}\), \(2 - x \to 0^{+}\), so
\[
\lim_{x \to 2^{-}} \sqrt{2 - x} = 0.
\]
As \(x \to 2^{+}\), \(2 - x \to 0^{-}\). The square root (even index) of a negative quantity is undefined in \(\R\), so the right-hand limit does not exist:
\[
\lim_{x \to 2^{+}} \sqrt{2 - x} \quad \text{dne}.
\]
\end{examplebox}


\begin{examplebox}{Example 7}

Consider \(g(x) = \sqrt{2x^{2} + x - 1} = \sqrt{(2x - 1)(x + 1)}\). Analyze the sign of the radicand near the critical points \(x = -1\) and \(x = \frac{1}{2}\):
\begin{align*}
\lim_{x \to -1^{-}} \sqrt{(2x - 1)(x + 1)}
    &= \sqrt{ ( -3\,)(\,0^{-} ) } = \sqrt{0^{+}} = 0,\\[6pt]
\lim_{x \to -1^{+}} \sqrt{(2x - 1)(x + 1)}
    &= \sqrt{ ( -3\,)(\,0^{+} ) } = \sqrt{0^{-}} \quad \text{dne},\\[6pt]
\lim_{x \to \frac{1}{2}^{-}} \sqrt{(2x - 1)(x + 1)}
    &= \sqrt{ (\,0^{-} )\bigl(\tfrac{3}{2}\bigr) } = \sqrt{0^{-}} \quad \text{dne},\\[6pt]
\lim_{x \to \frac{1}{2}^{+}} \sqrt{(2x - 1)(x + 1)}
    &= \sqrt{ (\,0^{+} )\bigl(\tfrac{3}{2}\bigr) } = \sqrt{0^{+}} = 0.
\end{align*}
\end{examplebox}


\begin{examplebox}{Example 8}

Let
\[
f(x) = \begin{cases}
\dfrac{x^{2} - x - 2}{x + 1}, & x < 0, \\[10pt]
\sqrt{4 - x},                  & 0 \le x \le 4, \\[6pt]
x^{2} - 5x + 4,                & x > 4.
\end{cases}
\]
Find each limit.
\begin{enumerate}
    \item[(a)] \(\displaystyle\lim_{x \to -1} f(x)\).
    For \(x\) near \(-1\) (but \(x \neq -1\)) we use the first piece. Factor the numerator:
    \[
    \frac{x^{2} - x - 2}{x + 1} = \frac{(x - 2)(x + 1)}{x + 1} = x - 2.
    \]
    Hence \(\displaystyle\lim_{x \to -1} f(x) = \lim_{x \to -1} (x - 2) = -3\).

    \item[(b)] \(\displaystyle\lim_{x \to 0^{-}} f(x)\).
    For \(x\) just below 0, the first piece applies:
    \[
    \lim_{x \to 0^{-}} \frac{x^{2} - x - 2}{x + 1} = \frac{0 - 0 - 2}{0 + 1} = -2.
    \]

    \item[(c)] \(\displaystyle\lim_{x \to 0^{+}} f(x)\).
    For \(x\) just above 0, use \(\sqrt{4 - x}\):
    \[
    \lim_{x \to 0^{+}} \sqrt{4 - x} = \sqrt{4} = 2.
    \]

    \item[(d)] \(\displaystyle\lim_{x \to 4^{+}} f(x)\).
    For \(x\) just above 4, the third piece applies:
    \[
    \lim_{x \to 4^{+}} (x^{2} - 5x + 4) = 16 - 20 + 4 = 0.
    \]
\end{enumerate}
\end{examplebox}

\bigskip


\subsection*{Two-Sided Limits via One-Sided Limits}
\label{subsec:1.2.2}


Some limits are most easily evaluated by first computing the two one-sided limits. The following theorem makes this connection precise.


\begin{theorembox}{Theorem 9}

\[
\lim_{x \to a} f(x) = L \quad\iff\quad \lim_{x \to a^{-}} f(x) = L = \lim_{x \to a^{+}} f(x).
\]
The two-sided limit exists if and only if both one-sided limits exist and are equal. When they are equal, the two-sided limit is that common value.
\end{theorembox}


\begin{examplebox}{Example 10}

From Example~6, \(\displaystyle\lim_{x \to 2^{-}} \sqrt{2 - x} = 0\) but \(\displaystyle\lim_{x \to 2^{+}} \sqrt{2 - x}\) does not exist. By Theorem~9, the two-sided limit cannot exist:
\[
\lim_{x \to 2} \sqrt{2 - x} \quad \text{dne}.
\]
\end{examplebox}


\begin{examplebox}{Example 11}

For the Heaviside function \(H(x)\) from Example~2, we have \(\lim_{x \to 0^{-}} H(x) = 0\) and \(\lim_{x \to 0^{+}} H(x) = 1\). Since the one-sided limits differ, Theorem~9 gives
\[
\lim_{x \to 0} H(x) \quad \text{dne}.
\]
\end{examplebox}


\begin{examplebox}{Example 12}

Let
\[
f(x) = \begin{cases}
\sqrt{4 - x},   & x \le 4, \\[4pt]
x^{2} - 5x + 4, & x > 4.
\end{cases}
\]
Evaluate \(\displaystyle\lim_{x \to 4} f(x)\).

\medskip
\textbf{Solution.}
The function is defined differently for \(x < 4\) and \(x > 4\), so we check one-sided limits.
\[
\lim_{x \to 4^{-}} f(x) = \lim_{x \to 4^{-}} \sqrt{4 - x} = 0,
\qquad
\lim_{x \to 4^{+}} f(x) = \lim_{x \to 4^{+}} (x^{2} - 5x + 4) = 0.
\]
Both one-sided limits equal 0, so by Theorem~9, \(\displaystyle\lim_{x \to 4} f(x) = 0\).
\end{examplebox}


\begin{examplebox}{Example 13}

Let \(g(t) = \dfrac{t + 4}{|t + 4|}\). Evaluate \(\displaystyle\lim_{t \to -4} g(t)\).

\medskip
\textbf{Solution.}
Recall the piecewise definition of the absolute value:
\[
|t + 4| = \begin{cases}
\phantom{-}t + 4,  & t \ge -4, \\[4pt]
-(t + 4),          & t < -4.
\end{cases}
\]
For the left-hand limit, \(t < -4\):
\[
\lim_{t \to -4^{-}} g(t) = \lim_{t \to -4^{-}} \frac{t + 4}{-(t + 4)} = \lim_{t \to -4^{-}} (-1) = -1.
\]
For the right-hand limit, \(t > -4\):
\[
\lim_{t \to -4^{+}} g(t) = \lim_{t \to -4^{+}} \frac{t + 4}{t + 4} = \lim_{t \to -4^{+}} 1 = 1.
\]
The one-sided limits differ, so \(\displaystyle\lim_{t \to -4} g(t)\) does not exist.
\end{examplebox}

\bigskip


\subsection*{Limits Involving the Greatest Integer Function}
\label{subsec:1.2.3}


The greatest integer function (or floor function), denoted \([\![x]\!]\), returns the largest integer less than or equal to \(x\). Its piecewise nature makes one-sided limits essential for evaluating limits at integer points.


\begin{examplebox}{Example 14}

Evaluate \(\displaystyle\lim_{x \to -2} [\![x]\!]\).

\medskip
\textbf{Solution.}
For values of \(x\) satisfying \(-3 \le x < -2\), we have \([\![x]\!] = -3\). For \(-2 \le x < -1\), we have \([\![x]\!] = -2\).
\begin{align*}
\lim_{x \to -2^{-}} [\![x]\!] &= \lim_{x \to -2^{-}} (-3) = -3,\\[6pt]
\lim_{x \to -2^{+}} [\![x]\!] &= \lim_{x \to -2^{+}} (-2) = -2.
\end{align*}
Since the one-sided limits differ, Theorem~9 gives \(\displaystyle\lim_{x \to -2} [\![x]\!]\) dne.
\end{examplebox}


\begin{examplebox}{Example 15}

Evaluate \(\displaystyle\lim_{s \to 2^{+}} \bigl(s + [\![1 - s]\!]\bigr)\).

\medskip
\textbf{Solution.}
For \(s\) just greater than 2, the quantity \(1 - s\) is just less than \(-1\). Since \(-2 \le 1 - s < -1\) for such \(s\), we have \([\![1 - s]\!] = -2\), and
\[
\lim_{s \to 2^{+}} \bigl(s + [\![1 - s]\!]\bigr)
    = \lim_{s \to 2^{+}} (s - 2) = 0.
\]
\end{examplebox}


\begin{examplebox}{Example 16}

Evaluate \(\displaystyle\lim_{x \to \frac{1}{2}^{-}} \frac{[\![2x - 1]\!] - 2x}{2x + 1}\).

\medskip
\textbf{Solution.}
As \(x \to \frac{1}{2}^{-}\), we have \(2x - 1 \to 0^{-}\). For \(x\) just below \(\frac{1}{2}\), the quantity \(2x - 1\) satisfies \(-1 \le 2x - 1 < 0\), so \([\![2x - 1]\!] = -1\).
\[
\frac{[\![2x - 1]\!] - 2x}{2x + 1}
    = \frac{-1 - 2x}{2x + 1}
    = -\frac{2x + 1}{2x + 1} = -1.
\]
The expression is identically \(-1\) for all \(x\) sufficiently close to \(\frac{1}{2}\) from the left, so
\[
\lim_{x \to \frac{1}{2}^{-}} \frac{[\![2x - 1]\!] - 2x}{2x + 1} = -1.
\]
\end{examplebox}


\begin{examplebox}{Example 17}

Evaluate \(\displaystyle\lim_{x \to 1} \frac{x^{2} - 2x + 1}{x - [\![x]\!]}\).

\medskip
\textbf{Solution.}
We inspect each one-sided limit and apply Theorem~9.

\textbf{Left-hand limit} (\(x \to 1^{-}\)): For \(0 \le x < 1\), \([\![x]\!] = 0\).
\[
\lim_{x \to 1^{-}} \frac{x^{2} - 2x + 1}{x - [\![x]\!]}
    = \lim_{x \to 1^{-}} \frac{(x - 1)^{2}}{x}
    = \frac{0}{1} = 0.
\]

\textbf{Right-hand limit} (\(x \to 1^{+}\)): For \(1 \le x < 2\), \([\![x]\!] = 1\).
\[
\lim_{x \to 1^{+}} \frac{x^{2} - 2x + 1}{x - [\![x]\!]}
    = \lim_{x \to 1^{+}} \frac{(x - 1)^{2}}{x - 1}
    = \lim_{x \to 1^{+}} (x - 1) = 0.
\]

Both one-sided limits equal 0, so \(\displaystyle\lim_{x \to 1} \frac{x^{2} - 2x + 1}{x - [\![x]\!]} = 0\).
\end{examplebox}

\pagebreak


\subsection*{Exercises}
\label{subsec:1.2.4}


\raggedcolumns
\setlength{\columnsep}{1.5em}

\medskip
\noindent\textbf{A.} Evaluate the following limits.

\begin{multicols}{2}
\begin{enumerate}[label=\arabic*., nosep, topsep=0pt, itemsep=2pt]
    \item \(\displaystyle\lim_{x \to 2} \sqrt{7x^{2} - 12x + 5}\)
    \item \(\displaystyle\lim_{x \to \frac{5}{2}} \sqrt{4x^{2} - 20x + 25}\)
    \item \(\displaystyle\lim_{x \to \frac{1}{2}^{+}} [\![2x + 1]\!]\)
    \item \(\displaystyle\lim_{x \to \frac{1}{2}^{-}} [\![2x + 1]\!]\)
    \item \(\displaystyle\lim_{x \to \frac{1}{2}} [\![2x + 1]\!]\)
    \item \(\displaystyle\lim_{x \to 2} \frac{|3x|}{|x - 4|\,(x - 2)}\)
    \item \(\displaystyle\lim_{t \to 1} \frac{3t - 2}{|2t^{2} - 5t + 3|}\)
    \item \(\displaystyle\lim_{x \to 0} \frac{x - [\![x]\!] - 1}{[\![x]\!] + 1}\)
\end{enumerate}
\end{multicols}

\medskip
\noindent\textbf{B.} Let
\[
g(x) = \begin{cases}
\dfrac{x^{2} - 9}{x^{2} + 3x}, & x \le -2, \\[10pt]
\sqrt{3x + 7} - 1,               & x > -2.
\end{cases}
\]
Evaluate: \(\displaystyle\lim_{x \to -3} g(x)\), \quad \(\displaystyle\lim_{x \to -2} g(x)\), \quad \(\displaystyle\lim_{x \to 1} g(x)\).

\medskip
\noindent\textbf{C.} Let
\[
f(x) = \begin{cases}
\dfrac{x^{2} + 5x + 6}{x^{2} - x - 6}, & x \le -2, \\[10pt]
[\![x + 2]\!],                           & -2 < x < 0, \\[6pt]
\sqrt{x + 1},                            & x \ge 0.
\end{cases}
\]
Evaluate: \(\displaystyle\lim_{x \to -2} f(x)\), \quad \(\displaystyle\lim_{x \to -1} f(x)\), \quad \(\displaystyle\lim_{x \to 0} f(x)\), \quad \(\displaystyle\lim_{x \to 3} f(x)\).

\medskip
\noindent\textbf{D.} Given
\[
h(x) = \begin{cases}
kx + 2,   & x \le 2, \\[4pt]
x^{2} - k, & x > 2,
\end{cases}
\]
where \(k\) is a constant. Find \(k\) so that \(\displaystyle\lim_{x \to 2} h(x)\) exists.

\medskip
\noindent\textbf{E.} Sketch a graph of a function \(f(x)\) satisfying all of the following:
\begin{multicols}{2}
\begin{itemize}[nosep, topsep=0pt, itemsep=3pt]
    \item \(\operatorname{dom} f = [-4, 4]\)
    \item \(f(-4) = f(-2) = 2\)
    \item \(f(0) = 1\)
    \item \(f(2) = -1\)
    \item \(f(4) = 0\)
    
    \columnbreak
    
    \item \(\displaystyle\lim_{x \to -4^{+}} f(x) = 0\)
    \item \(\displaystyle\lim_{x \to -2} f(x) = 3\)
    \item \(\displaystyle\lim_{x \to 0^{-}} f(x) = -1\)
    \item \(\displaystyle\lim_{x \to 0^{+}} f(x) = 4\)
    \item \(\displaystyle\lim_{x \to 2} f(x) = -1\)
    \item \(\displaystyle\lim_{x \to 4} f(x) = 0\)
\end{itemize}
\end{multicols}
